\begin{pastexam}[label={px:ncu-113-dsa-15}]{中央資工所 113 學年度 第 15 題}
（複選；每題 5 分，全部答對才給分，答錯倒扣一分）

The following statements about minimum spanning tree (MST) may or may not be correct. Assume that the weighted
graph $G=(V,E)$ is undirected and connected. Do not assume that edge weights are distinct unless this is
specifically stated. Chose the correct items.
\begin{enumerate}[label=(\Alph*)]
\item If $|E|\ge|V|-1$, and there is a unique heaviest edge, then this edge cannot be part of any MST of $G$.
\item If $G$ has a cycle with a unique heaviest edge $e$, then $e$ cannot be part of any MST of $G$.
\item If the lightest edge in $G$ is unique, then it must be part of every MST of $G$.
\item If $G$ has a cycle with a unique lightest edge $e$, then $e$ must be part of some MST of $G$.
\item The shortest path between two nodes is necessarily part of some MST of $G$.
\end{enumerate}
\end{pastexam}

\begin{pastsoln}
答案：(B)(C)。

\begin{itemize}
\item (A) 錯。 最重的邊如果是橋（拿掉就不連通），每棵生成樹都要用它。 反例：$G$ 本身就是一棵樹，三個點
$a$--$b$ 權重 $1$、$b$--$c$ 權重 $5$，$|E| = 2 = |V| - 1$，最重的邊 $b$--$c$ 唯一，卻在唯一的 MST 裡。
\item (B) 對。 這是環性質（資結書定理~24.3）。 證明：設某棵 MST $T$ 含 $e = (u, v)$。 從 $T$ 拿掉 $e$，
$T$ 分成兩塊，$u$、$v$ 各在一塊。 環上除了 $e$ 以外的部分是一條 $u$ 到 $v$ 的路，一定有某條邊 $f$ 跨越這
兩塊。 $w(f) < w(e)$（$e$ 是環上唯一最重的），$T - e + f$ 是生成樹而且更輕，和 $T$ 是最小的矛盾。
\item (C) 對。 設 $e = (u, v)$ 是唯一最輕的邊，某棵 MST $T$ 不含 $e$。 把 $e$ 加進 $T$ 形成一個環，環上
其他邊都比 $e$ 重，拿掉其中任一條 $f$，$T + e - f$ 是更輕的生成樹，矛盾。 所以每棵 MST 都含 $e$。
\item (D) 錯。 $e$ 在某個環上最輕，不代表它在別的環上不是最重。 反例：四個點 $u, v, x, y$，邊
$u$--$v$ 權重 $5$、$v$--$x$ 權重 $10$、$x$--$y$ 權重 $10$、$y$--$u$ 權重 $10$，再加一個點 $z$ 和兩條邊
$u$--$z$、$z$--$v$，權重都是 $1$。 $e = u$--$v$ 在環 $u, v, x, y$ 上是唯一最輕的；但在環 $u, z, v$ 上它是
唯一最重的，由 (B)，它不在任何 MST 裡。
\item (E) 錯。 三角形 $a$--$b$ 權重 $3$、$b$--$c$ 權重 $3$、$a$--$c$ 權重 $4$。 $a$ 到 $c$ 的最短路徑是直接
走 $a$--$c$（$4 < 6$），但 $a$--$c$ 是這個環上唯一最重的邊，不在任何 MST 裡（唯一的 MST 是 $a$--$b$、
$b$--$c$）。 最小生成樹讓「總權重」最小，不保證兩點之間的路徑最短（資結書第~24.5.4~節）。
\end{itemize}
\end{pastsoln}
